% AY2020 制御工学 第2問〔II〕（転記）
% 出典: 03_制御工学/original/03_制御工学/2020/2020se.pdf
% 要約: G1,G2 直列の前向き経路に K の負帰還を持つフィードバック系。
%       一巡伝達関数のゲイン位相・ステップ応答、K 調整、G2 変更時の比較考察。

\section*{第2問〔II〕}

制御システムに関する次の問い (1)〜(4) に答えよ. なお, 伝達関数は $G_1(s)=\dfrac{K_1}{s+1}$,
$G_2(s)=\dfrac{1}{0.1s+1}$ とする.

\begin{center}
\begin{tikzpicture}[>=Latex]
  \node (in)  at (0,0) {};
  \node (sum) [sum] at (1.0,0) {};
  \node (G1)  [block] at (2.6,0) {$G_1(s)$};
  \node (G2)  [block] at (4.6,0) {$G_2(s)$};
  \coordinate (nd) at (5.9,0);
  \node (out) at (7.2,0) {};
  \node (K)   [block] at (3.6,-1.5) {$K$};

  \draw[->] (in)  -- (sum);
  \draw[->] (sum) -- (G1);
  \draw[->] (G1)  -- (G2);
  \draw[->] (G2)  -- (out);
  \fill (nd) circle (1.4pt);
  \draw (nd) -- (5.9,-1.5) -- (K.east);
  \draw[->] (K.west) -| (sum);

  \node at ($(sum.north west)+(0.02,0.10)$) {\scriptsize $+$};
  \node at ($(sum.south west)+(0.02,-0.10)$) {\scriptsize $-$};
  \node at (3.6,-2.2) {\small 図2};
\end{tikzpicture}
\end{center}

\begin{enumerate}[label=(\arabic*)]
  \item 図2のシステムについて, 一巡伝達関数のゲインと位相を求めよ. ここで, $K_1=2$ とし,
        フィードバックゲインは $K=1$ とする. 加えて, フィードバックシステムに単位ステップ
        入力を与えた時の応答を求めよ.

  \item 一巡伝達関数のゲイン線図を漸近線で描け. 折れ点等の特徴点を明示し,
        解答用紙内の縦軸・横軸の目盛には適宜数値を記入すること.

  \item フィードバックゲイン $K$ を調整し, $\omega=0.01$ 近傍のゲインをおおよそ $0$ dB と
        したい. このときに必要なフィードバックゲイン $K$ の値を求めよ.

  \item 次に, 問い (3) で求めたフィードバックゲイン $K$ を用いて, $G_2(s)=\dfrac{1}{0.2s+1}$
        にした場合を考える. まず, 一巡伝達関数のゲイン線図を問い (2) で描いたゲイン線図上に
        破線で描き, 単位ステップ入力を与えた時の応答を求めよ. 次に, 問い (2) と (3) で
        得られたゲイン線図や応答と比較し, 2 つのシステムの違いについて具体的に考察せよ.
\end{enumerate}
